\chapter{\texttt{pipeline.solar} --- solar production to surface}
\label{ch:pipeline-solar}

This part covers the five \texttt{pipeline.*} workflows: user-facing
orchestration layers sitting above \texttt{core} and \texttt{medium}
(\cref{ch:core-common,ch:core-bsm,ch:core-numerical,ch:core-perturbative,%
ch:vacuum,ch:solar,ch:earth,ch:atmosphere}). None of the five chapters that
follow introduce new physics: every pipeline function's job is to prepare
inputs, call the physics blocks already derived in the previous parts of
this document, integrate or aggregate the resulting observables, and return
them bundled in a single result object -- so, unlike earlier chapters, none
of these sections carry a blue \textbf{Theoretical Foundation} box except
where a pipeline composes two physics stages in a way that is not already
written down elsewhere. Each chapter instead closes with a teal
\textbf{Usage Example} box: a complete, runnable Python snippet showing how
that pipeline's entry points are actually called together.

\section{\texttt{config.propagation.PropagationConfig} --- shared pipeline configuration}
\label{sec:propagation-config}

\modulehead{config/propagation.py}{\pyclass{PropagationConfig}: the single
configuration object every pipeline function in this part accepts.}

\begin{implbox}
\pyclass{PropagationConfig} replaces what used to be 20--38 separate
keyword arguments threaded into each pipeline entry point with one
dataclass. There is exactly one \pyclass{PropagationConfig} for the whole
project, not one per pipeline: fields a given pipeline does not need (e.g.\
\code{exposure}/\code{solar} for the atmosphere pipeline,
\cref{ch:pipeline-atmosphere}) are simply left at their defaults and
unused. Pipeline functions destructure it into the smaller pieces
\texttt{core}/\texttt{medium} functions actually need (a
\pyclass{RuntimeContext}, an \pyclass{OscillationParameters}, a built
profile object, \ldots) rather than passing the whole object down, so those
lower layers stay decoupled from this pipeline-level container.

\begin{center}
\small
\begin{tabular}{@{}llll@{}}
\toprule
Field & Type & Default & Role \\
\midrule
\code{runtime} & \pyclass{RuntimeContext} & required &
  \parbox[t]{4.6cm}{Torch device/dtype used throughout.} \\[4pt]
\code{oscillation} & \pyclass{OscillationParameters} & required &
  \parbox[t]{4.6cm}{Built \code{pmns} object plus mass splittings and
  antinu selection (\cref{ch:core-common}); see
  \pyfunc{oscillation\_parameters\_from\_preset} below.} \\[8pt]
\code{exposure} & \pyclass{ExposureParameters} & \code{ExposureParameters()} &
  \parbox[t]{4.6cm}{Nadir-exposure table settings (\cref{sec:earth-exposure}).
  Used by \texttt{pipeline.earth}/\texttt{solar\_earth} when averaging over a
  detector exposure window; unused by atmosphere.} \\[10pt]
\code{earth} & \pyclass{EarthParameters} & \code{EarthParameters()} &
  \parbox[t]{4.6cm}{Earth density-profile construction settings
  (\cref{sec:earth-profile}). Used by every pipeline with an
  Earth-crossing stage.} \\[8pt]
\code{solar} & \pyclass{SolarParameters} & \code{SolarParameters()} &
  \parbox[t]{4.6cm}{Solar-model construction settings (\cref{ch:solar}).
  Solar pipelines only.} \\[6pt]
\code{atmosphere} & \pyclass{AtmosphereParameters} & \code{AtmosphereParameters()} &
  \parbox[t]{4.6cm}{Atmosphere density-profile construction settings
  (\cref{ch:atmosphere}). Atmosphere pipeline only.} \\[6pt]
\code{production\_mode} & \code{"point"/"coherent"/"incoherent"} & \code{"point"} &
  \parbox[t]{4.6cm}{Solar production-radius treatment (solar pipelines
  only).} \\[6pt]
\code{rho0} & \code{Optional[TensorLike]} & \code{None} &
  \parbox[t]{4.6cm}{Production radius (fraction of $R_\odot$) for
  \code{production\_mode="point"}.} \\[6pt]
\code{source} & \code{Optional[str]} & \code{None} &
  \parbox[t]{4.6cm}{Solar source key (e.g.\ \code{"8B"}, \code{"pp"}) for
  \code{production\_mode} in \{\code{"coherent"}, \code{"incoherent"}\}.} \\[6pt]
\code{detector\_depth\_m} & \code{float} & \code{0.0} &
  \parbox[t]{4.6cm}{Detector depth below the Earth's surface, in metres;
  shifts the Earth-crossing path length wherever an Earth-regeneration
  stage is present.} \\[8pt]
\code{reunitarize\_earth} & \code{bool} & \code{False} &
  \parbox[t]{4.6cm}{Project the Earth evolution operator back onto the
  nearest unitary matrix, absorbing small numerical drift.} \\[6pt]
\code{reunitarize\_atmosphere} & \code{bool} & \code{False} &
  \parbox[t]{4.6cm}{Same, for the atmosphere evolution operator
  (\pyfunc{atmosphere\_evolutor}). Independent of
  \code{reunitarize\_earth}: the two stages are unitarized separately.} \\[8pt]
\code{analytic\_eigenvalues} & \code{bool} & \code{False} &
  \parbox[t]{4.6cm}{Use the closed-form Cardano/Ferrari eigenvalue solution
  (\cref{ch:core-perturbative}) instead of \code{torch.linalg.eigvalsh} for
  \code{method="analytical"} Earth/atmosphere propagation; raises if
  combined with \code{method="numerical"}, which never diagonalises.} \\[10pt]
\bottomrule
\end{tabular}
\end{center}

\pyfunc{PropagationConfig.oscillation\_parameters\_from\_preset(preset\_name="\_SM\_NUFIT52\_NO", *, antinu=False, NSI\_extension=None, context=None)}
is the static factory that builds the \code{oscillation} field
(\cref{ch:core-common}, Default Configuration box): it looks up
\code{preset\_name} in \code{tpeanuts.config.presets.OSCILLATION\_PRESETS},
dispatches to a sterile or plain-SM \pyclass{PMNS} exactly as described in
\cref{ch:core-bsm}, and optionally attaches an \pyclass{NSIConfig} built
from \code{NSI\_extension} (\cref{sec:nsi-presets}). This is the one place
in the pipeline layer allowed to depend on concrete SM/BSM/NSI/preset
implementations, keeping \pyclass{OscillationParameters} itself a passive
container (\cref{ch:core-common}).
\end{implbox}

\section{\texttt{pipeline.solar} --- solar production to surface}

\modulehead{pipeline/solar.py}{\pyfunc{propagate\_solar\_to\_surface}: the
incoherent solar mass mixture produced inside the Sun, at the solar
surface.}

\begin{implbox}
\begin{itemize}
  \item \pyclass{SolarSurfaceResult} (frozen dataclass): \code{source},
    \code{E\_MeV}, the built \code{profile}
    (\pyclass{medium.solar.profile.SolarProfile}), the source's
    \code{production\_distribution} (production-radius weights), the
    incoherent \code{mass\_weights} $w_i(E)$, and the corresponding
    \code{flavour\_probabilities} -- the same quantities
    \pyfunc{solar\_probability\_mass}/\pyfunc{probability\_incoherent}
    (\cref{ch:solar,ch:core-common}) already compute, simply bundled
    together with the profile and inputs that produced them.
  \item \pyfunc{propagate\_solar\_to\_surface(*, E\_MeV, config, source,
    solar\_profile=None, method="adiabatic\_approximated",
    legacy\_precision=False, include\_matter\_nc=None)}: builds (or reuses
    \code{solar\_profile}), reads off the source's production-radius
    distribution, calls \pyfunc{solar\_probability\_mass} for the mass
    weights, then \pyfunc{probability\_incoherent} with the identity
    operator standing in for ``no further coherent evolution'' to obtain
    \code{flavour\_probabilities} at the solar surface directly (before any
    Earth regeneration). \code{method} and \code{include\_matter\_nc} are
    exactly the ones documented in \cref{ch:solar}'s
    \texttt{probability.py} section; this function adds no behaviour of
    its own beyond wiring \pyclass{PropagationConfig} into that call.
\end{itemize}
\end{implbox}

\begin{examplebox}[title={Solar mass mixture at two energies, boron-8 source}]
\begin{lstlisting}[style=tpython]
import torch

from tpeanuts.config.propagation import PropagationConfig
from tpeanuts.pipeline.solar import propagate_solar_to_surface
from tpeanuts.util.context import RuntimeContext

context = RuntimeContext.resolve("cpu", torch.float64)
oscillation = PropagationConfig.oscillation_parameters_from_preset(
    "_SM_NUFIT52_NO", context=context,
)
config = PropagationConfig(runtime=context, oscillation=oscillation)

result = propagate_solar_to_surface(
    E_MeV=[2.0, 8.0],
    config=config,
    source="8B",
)

print(result.mass_weights.shape)          # torch.Size([2, 3])
print(result.mass_weights.sum(dim=-1))    # tensor([1., 1.], dtype=torch.float64)
print(result.flavour_probabilities)       # P(nu_e), P(nu_mu), P(nu_tau) per energy
\end{lstlisting}
\end{examplebox}

The default \code{method="adiabatic\_approximated"} uses the closed-form
two-level mixing-angle reduction (\cref{ch:solar}); passing
\code{method="adiabatic\_exact"} or \code{method="numerical"} instead --
required for NSI or a sterile \code{oscillation} preset, per the support
matrix in \cref{sec:bsm-support-matrix} -- needs no other change to the
call above.
